\section{Kiến trúc và Phương pháp đề xuất}
Nghiên cứu đề xuất kiến trúc hai giai đoạn: (1) Tiền xử lý dữ liệu với lấy mẫu lai RUS-ENN và huấn luyện mô hình Giáo viên (Stacking Teacher); (2) Áp dụng chưng cất tri thức (Knowledge Distillation) để chuyển giao tri thức từ Teacher sang mô hình Học viên (Student) Cây quyết định.

\subsection{Xử lý nhiễu và phân phối dữ liệu}
\subsubsection{Làm sạch và chuẩn hóa đặc trưng} 
Dữ liệu CICIDS2017 được làm sạch (loại bỏ \texttt{NaN}, \texttt{Inf}, trùng lặp) và chuẩn hóa bằng Robust Scaling để giảm thiểu ảnh hưởng của ngoại lai:
\begin{equation}
x_{norm} = \frac{x_i - \tilde{x}}{\text{IQR}(X)}
\end{equation}
trong đó $\tilde{x}$ là trung vị và $\text{IQR}(X)$ là khoảng tứ phân vị.

\subsubsection{Chiến lược lấy mẫu lai (RUS kết hợp ENN)}
Để giải quyết mất cân bằng lớn (lớp thiểu số $<0,01\%$), chúng tôi áp dụng chiến lược lấy mẫu lai. Đầu tiên, kỹ thuật Lấy mẫu giảm ngẫu nhiên (Random Under-Sampling - RUS) được sử dụng để giảm quy mô lớp đa số, loại bỏ bớt các mẫu cách xa ranh giới quyết định. Kế tiếp, bộ lọc ENN~\cite{wilson1972asymptotic} loại bỏ các mẫu bị nhiễu cục bộ dựa trên quy tắc 3 láng giềng gần nhất. Quy trình này (Thuật toán~\ref{algo:preprocessing}) tạo ra tập dữ liệu cân bằng với ranh giới phân lớp sắc nét.

\begin{algorithm}[hbt!]
\SetAlgoLined
\KwIn{Bộ dữ liệu thô $\mathcal{D} = \mathcal{S}_{maj} \cup \mathcal{S}_{min}$}
\KwOut{Bộ dữ liệu cân bằng $\mathcal{D}_{bal}$}
Loại bỏ \texttt{NaN}, \texttt{Inf}, trùng lặp và chuẩn hóa Robust Scaler\;
Giảm lấy mẫu $\mathcal{S}_{maj}$ bằng RUS để tạo $\mathcal{S}'_{maj}$\;
Khởi tạo $\mathcal{D}_{bal} \leftarrow \mathcal{S}'_{maj} \cup \mathcal{S}_{min}$\;
\ForEach{$(x_i, y_i) \in \mathcal{S}'_{maj}$}{
    Tìm 3 láng giềng gần nhất $N_3(x_i)$ trong $\mathcal{D}_{bal}$\;
    \If{nhãn đa số của $N_3(x_i) \neq y_i$}{
        $\mathcal{D}_{bal} \leftarrow \mathcal{D}_{bal} \setminus \{(x_i, y_i)\}$\;
    }
}
\Return{$\mathcal{D}_{bal}$}
\caption{Tiền xử lý và Lấy mẫu lai (RUS + ENN)}
\label{algo:preprocessing}
\end{algorithm}

\subsection{Mô hình Teacher kiến trúc Stacking}
Kiến trúc Stacking gồm hai tầng:
\subsubsection{Tầng 0 (Nhóm mô hình cơ sở)} Bao gồm \textbf{XGBoost}~\cite{chen2016xgboost}, \textbf{LightGBM}~\cite{ke2017lightgbm} (cơ chế boosting) và \textbf{Random Forest}~\cite{breiman2001random} (cơ chế bagging).
\subsubsection{Tầng 1 (Mô hình tổng hợp)} Đầu ra xác suất từ Tầng 0 tạo thành ma trận đặc trưng $Z_i = [P_{XGB}, P_{LGBM}, P_{RF}]$. Bộ phân loại Hồi quy Logistic tổng hợp dự đoán qua hàm sigmoid: $\hat{Y} = \sigma(W \cdot Z_i + b)$. Kiến trúc này tăng độ chính xác và giảm cảnh báo giả.

\subsection{Chưng cất tri thức sang mô hình Student}
Mô hình Stacking Teacher ($\mathcal{E}$) có độ trễ lớn, không phù hợp cho thiết bị biên. Do đó, chúng tôi áp dụng chưng cất tri thức dạng Mô hình thay thế (Surrogate Model) để huấn luyện một mô hình Student (Cây quyết định nông). Cây quyết định được chọn vì tính gọn nhẹ và tốc độ duyệt $O(\log n)$.

Cụ thể, mô hình Teacher tổng hợp nhãn mục tiêu mới $\hat{Y}_{teacher}$ bằng cách tạo ra các dự đoán nhãn cứng (hard labels) trên tập huấn luyện. Sau đó, mô hình Student $M_{student}$ học trực tiếp trên bộ nhãn này (Thuật toán~\ref{algo:kd}).

\begin{algorithm}[ht]
\DontPrintSemicolon
\SetAlgoLined
\KwIn{Tập huấn luyện $X_{train} = \{x_1, \dots, x_N\}$, Mô hình Giáo viên $\mathcal{E}$}
\KwOut{Mô hình Học viên $M_{student}$}
Khởi tạo mảng dự đoán $\hat{Y}_{teacher} \leftarrow$ mảng rỗng kích thước $N$\;
\For{$i = 1 \dots N$}{
    $p_i \leftarrow \mathcal{E}.predict\_proba(x_i)$\;
    $\hat{Y}_{teacher}[i] \leftarrow \arg\max(p_i)$ \;
}
Khởi tạo Cây quyết định $M_{student}$ (max\_depth=$d$)\;
$M_{student}.fit(X_{train}, \hat{Y}_{teacher})$ \;
\Return $M_{student}$ \;
\caption{Chưng cất tri thức Giáo viên - Học viên}
\label{algo:kd}
\end{algorithm}
